The basic equation of radioactive decay can be expressed as:
\[ N(t) = N_0 \exp(-\lambda t) \]
where $N(t)$ and $N_0$ are the number of remaining atoms of the radioactive
isotope (or parent isotope) at time $t$ and its initial number at $t = 0$,
respectively, while $\lambda$ is the decay constant. The decay of the parent
produces daughter nuclides $D(t)$, or radiogenics, which is defined as
\[ D(t) = N_0 - N(t)\,. \]

Based on those ideas, a group of astronomers investigates a number of
meteorite samples to determine their ages. They have two kinds of samples:
allende chondrite (A) and basaltic achondrite (B). From the samples, they
measure the abundance of $^{87}$Rb and $^{87}$Sr, where it is assumed that
$^{87}$Sr is entirely produced by the decay of $^{87}$Rb. The value of
$\lambda$ is $1.42 \times 10^{-11}$ per year for this isotopic decay. In
addition, non-radiogenic element $^{86}$Sr is also measured. Results of
measurement are given in the table below, expressed in ppm (part per
million).

\begin{center}
\begin{tabular}{c|c|c|c|c}
Sample No & Meteorite type & $^{86}$Sr (ppm) & $^{87}$Rb (ppm)
  & $^{87}$Sr (ppm) \\
\hline
1  & A & 29.6 & 0.3   & 20.7 \\
2  & B & 58.7 & 68.5  & 44.7 \\
3  & B & 74.2 & 14.4  & 52.9 \\
4  & A & 40.2 & 7.0   & 28.6 \\
5  & A & 19.7 & 0.4   & 13.8 \\
6  & B & 37.9 & 31.6  & 28.4 \\
7  & A & 33.4 & 4.0   & 23.6 \\
8  & B & 29.8 & 105.0 & 26.4 \\
9  & A & 9.8  & 0.8   & 6.9  \\
10 & B & 18.5 & 44.0  & 15.4 \\
\end{tabular}
\end{center}

Questions:
\begin{enumerate}
\item Express the time $t$ in term of $\dfrac{D(t)}{N(t)}$.
\item Determine the half-life $t_{1/2}$, i.e., the time needed to obtain a
  half number of parents after decay.
\item Knowledge on the ratio between two isotopes is more valuable than just
  the absolute abundance of each isotope. It is quite likely that there was
  some initial strontium present. By taking
  $\left(\frac{^{87}\mathrm{Rb}}{^{86}\mathrm{Sr}}\right)$ as independent
  variable and $\left(\frac{^{87}\mathrm{Sr}}{^{86}\mathrm{Sr}}\right)$ as
  dependent variable, estimate the simple linear regression model to
  represent the data.
\item Plot $\left(\frac{^{87}\mathrm{Rb}}{^{86}\mathrm{Sr}}\right)$ versus
  $\left(\frac{^{87}\mathrm{Sr}}{^{86}\mathrm{Sr}}\right)$ and also the
  regression line (isochrone) for each type of the meteorites. (Please use
  minimum 7 decimal digits for intermediate calculations)
\item Subsequently, help this astronomer to determine the age of each type of
  the meteorites and its error. Which type is older?
\item Determine the initial value of
  $\left(\frac{^{87}\mathrm{Sr}}{^{86}\mathrm{Sr}}\right)_0$ for each type of
  the meteorites and its error.
\end{enumerate}

\textit{Glossary}:

A simple linear regression line $y = a + bx$ can be fitted to a set of data
$(X_i, Y_i)$, $i = 1, \ldots, n$, in which
\[ b = \frac{SS_{xy}}{SS_{xx}} \]
\[ a = \bar{y} - b\bar{x} \]
where
\[ SS_{xx} \text{ : sum of square for } X
   = \sum_{i=1}^{n} X_i^2
   - \frac{1}{n}\left( \sum_{i=1}^{n} X_i \right)^2 \]
\[ SS_{yy} \text{ : sum of square for } Y
   = \sum_{i=1}^{n} Y_i^2
   - \frac{1}{n}\left( \sum_{i=1}^{n} Y_i \right)^2 \]
\[ SS_{xy} \text{ : sum of square for both } X \text{ and } Y
   = \sum_{i=1}^{n} X_i Y_i
   - \frac{1}{n} \sum_{i=1}^{n} X_i \sum_{i=1}^{n} Y_i \]

Standard deviation of each parameter, $a$ and $b$ can be calculated by
\[ S_a = \sqrt{ \frac{SS_{YY} - \dfrac{(SS_{XY})^2}{SS_{XX}}}
   {(n-2)\,SS_{XX}} \times \sum_{i=1}^{n} X_i^2 } \]
\[ S_b = \sqrt{ \frac{SS_{YY} - \dfrac{(SS_{XY})^2}{SS_{XX}}}
   {(n-2)\,SS_{XX}} } \]
